\section{Discussion}
\label{sec:discussion}

The most reliable improvement belongs to the evaluated VAE/GMM/IG package, not
to neighbourhood context alone. The centre-only result shows that much of the
gain over legacy msiPL remains when neighbours are removed. Context contributes
positively in the frozen CAC comparison, but the GBM, shuffled and attention
controls show that this effect cannot yet be generalised as a benefit caused by
correct local topology.

This distinction explains why a negative or mixed spatial result remains
useful. Uniform averaging compresses neighbouring spectra into one vector and
can mix tissue across boundaries. Learned attention can reweight neighbours but
was close to uniform in typical evaluated pixels and did not yield a stable
peak-quality advantage on the outcome-known pilot sections. These findings
bound the evaluated mechanisms; they do not imply that spatial modelling is
unhelpful in MSI or that convolutional, graph-based or multiscale approaches
would behave similarly.

The evaluation is spatial rather than molecular. A high PCC or mSCF1 indicates
that an ion image agrees with an expert mask under the chosen thresholds. It
does not identify the molecule, validate a biomarker or prove that a poorly
correlated ion is noise. Likewise, the GBM collection contains eight sections
from four patients. Patient-grouped descriptive summaries can show whether an
effect is driven by one patient, but the study does not support broad
population-level inference.

S3PL supplies two further cautions. Its CAC result demonstrates that its native
pipeline can exploit context effectively under the evaluated conditions, while
centre-tiled and topology controls answer narrower and partly confounded
questions. Its short GBM protocol is also sensitive to initialization. These
results make a single favourable score an inadequate baseline and show why
reconstruction loss alone cannot stand in for downstream peak quality.

